\clearpage
\VerdierPhysicalPageAnchor{0179}%
\begingroup\small\setlength{\parskip}{0pt}\setlength{\emergencystretch}{1.5em}
\setlength{\abovedisplayskip}{4pt plus 2pt minus 2pt}%
\setlength{\belowdisplayskip}{4pt plus 2pt minus 2pt}%

\VerdierSemanticAnchor{ch3d-1-2-4}
\noindent\textbf{1.2.4.}\quad
La sous-catégorie $K^{,b}(\mathcal A)/\operatorname{Ac}(\mathcal A)$
(resp. $K^{,+}(\mathcal A)/\operatorname{Ac}(\mathcal A)$, resp.
$K^{,-}(\mathcal A)/\operatorname{Ac}(\mathcal A)$) de $D(\mathcal A)$
est notée $D^b(\mathcal A)$ (resp. $D^+(\mathcal A)$, resp.
$D^-(\mathcal A)$) et est appelée la \emph{sous-catégorie des objets
bornés} (resp. \emph{bornés inférieurement}, resp. \emph{bornés
supérieurement}). Il convient de noter que ces limitations ne s'appliquent
qu'aux objets de cohomologie des complexes considérés.

\VerdierSemanticAnchor{ch3d-1-2-5}
\medskip\noindent\textbf{1.2.5.}\quad
Soient $X$ et $Y$ deux objets de $D(\mathcal A)$. On désigne par
$\operatorname{Ext}_{\mathcal A}^{0}(X,Y)$ le groupe
$\operatorname{Hom}_{D(\mathcal A)}(X,Y)$. De même, on désigne par
$\operatorname{Ext}_{\mathcal A}^{i}(X,Y)$ les groupes
$\operatorname{Ext}_{\mathcal A}^{0}(X,Y[i])\simeq
\operatorname{Ext}_{\mathcal A}^{0}(X[-i],Y)$ (chap. I, 1.3.8). Cette
notation sera justifiée ultérieurement. Rappelons cependant que tout triangle
distingué $(X,X',X'',u,v,w)$ (resp. $(Y,Y',Y'',u,v,w)$) de $D(\mathcal A)$
fournit une suite exacte illimitée :
\[
\cdots\longrightarrow \operatorname{Ext}_{\mathcal A}^{0}(X'',Y)
\longrightarrow \operatorname{Ext}_{\mathcal A}^{0}(X',Y)
\longrightarrow \operatorname{Ext}_{\mathcal A}^{0}(X,Y)
\longrightarrow \operatorname{Ext}_{\mathcal A}^{1}(X'',Y)
\longrightarrow\cdots
\]
(resp.
\[
\cdots\longrightarrow \operatorname{Ext}_{\mathcal A}^{0}(X,Y)
\longrightarrow \operatorname{Ext}_{\mathcal A}^{0}(X,Y')
\longrightarrow \operatorname{Ext}_{\mathcal A}^{0}(X,Y'')
\longrightarrow \operatorname{Ext}_{\mathcal A}^{1}(X,Y)
\longrightarrow\cdots )
\]
(chap. II, 1.2.1).

\VerdierSemanticAnchor{ch3d-1-2-6}
\medskip\noindent\textbf{1.2.6.}\quad
Soit $M$ un ensemble d'objets de $\mathcal A$. On désigne par
$D_M(\mathcal A)$ la sous-catégorie pleine de $D(\mathcal A)$ définie par
les complexes dont les objets de cohomologie appartiennent à $M$ en tout
degré. On désigne par $D_M^*(\mathcal A)$ la sous-catégorie pleine
$D_M(\mathcal A)\cap D^*(\mathcal A)$. Si la sous-catégorie pleine de
$\mathcal A$ définie par les objets de $M$ est une sous-catégorie épaisse de
$\mathcal A$ (\textit{i.e.} si elle est abélienne, si le foncteur
d'inclusion est exact et si l'ensemble des objets de $M$ est stable par
extension), les sous-catégories $D_M^*(\mathcal A)$ sont des
sous-catégories triangulées pleines de $D(\mathcal A)$.

\VerdierSemanticAnchor{ch3d-1-2-7}
\medskip\noindent\textbf{1.2.7.}\quad
Soient $F\colon\mathcal A\longrightarrow\mathcal B$ un foncteur exact entre
deux catégories abéliennes et
$K(F)\colon K(\mathcal A)\longrightarrow K(\mathcal B)$ l'extension du
foncteur $F$ aux complexes à homotopie près (chap. I, 2.5.10). Le foncteur
$K(F)$ transforme les complexes acycliques de $\mathcal A$ en complexes
acycliques de $\mathcal B$, est exact (chap. II, 1.3.7) et, par suite,
(chap. II, 2.2.6) définit un foncteur exact
$D(F)\colon D(\mathcal A)\longrightarrow D(\mathcal B)$. Le foncteur
$D(F)$ transforme les objets bornés (resp.\ bornés inférieurement, resp.\
bornés supérieurement) de $D(\mathcal A)$ en objets de même nature de
$D(\mathcal B)$. Le plus souvent, on utilisera la notation abusive
$F\colon D(\mathcal A)\longrightarrow D(\mathcal B)$ pour désigner le
foncteur $D(F)$.

\VerdierSemanticAnchor{ch3d-1-2-8}
\medskip\noindent\textbf{1.2.8.}\quad
L'isomorphisme canonique (chap. II, 1.3.8)
$i\colon K(\mathcal A)^o\longrightarrow K(\mathcal A^o)$ transforme les
objets de $\operatorname{Ac}(\mathcal A)^o$ en objets de
$\operatorname{Ac}(\mathcal A^o)$ et, par suite, définit un isomorphisme de
catégories triangulées encore noté
$i\colon D(\mathcal A)^o\longrightarrow D(\mathcal A^o)$.
\endgroup
